\newcommand{\Fdelta}{$\mathbb{F}$-Delta}
\newcommand{\PredErr}{Pred-Err}

As an illustration, consider a binomial model with parameter $\theta$.
Let $S_n = (X_1, X_2, \ldots, X_n)$ denote a sample of $n$ independent observations,
where $X_i \sim \operatorname{Bernoulli}(\theta)$.

We use the basis functions
\[
    b_k(s_n) \coloneqq \left(\frac1n \sum_{i=1}^n x_i \right)^k.
\]
These basis functions form polynomial features of the sufficient statistic.

For a binomial model with success probability $\theta$, the entropy of $S_n$ is
\[
    \entropy{Q}{S_n \mid \paramVal} = -n \left(\paramVal \ln \paramVal + (1 - \paramVal) \ln (1 - \paramVal)  \right).
\]

We consider the case $n=7$ and approximate $\cDist(S_7)$ using the procedure described above.
The parameter space is represented by a sample set \set{S} consisting of 30 values of $\theta$.

When the approximation lies close to a marginal likelihood within the binomial family,
the bound in Theorem~\ref{th:central-approx} can be estimated by
\[
    \max_{\inSet[\paramVal]{S}} \engFunc - \min_{\inSet[\paramVal]{S}} \engFunc,
\]
which we denote by \Fdelta.

To investigate the effect of the number of basis functions on the quality of the approximation, we vary the number
of basis functions from 4 to 14.
The corresponding values of \Fdelta\ are shown in Figure~\ref{fig:error_comparision}.

\begin{figure}[t]
    \centering
    \begin{tikzpicture}
        \begin{axis}[
            globalstyle,
            legend pos=north east,
            xlabel={Number of Basis},
            xmin=4, xmax=14,
            ymin=0.0
        %xtick={3,...,14},
        ]

            \addplot+[mark=*, blue] coordinates {
                (4, 0.622701) (5, 0.192069) (6, 0.192069) (7, 0.076361) (8, 0.076361) (9, 0.028878) (10, 0.028878) (11, 0.011228) (12, 0.015049) (13, 0.012935) (14, 0.014727)
            };
            \addlegendentry{\Fdelta}

        \end{axis}
    \end{tikzpicture}
    \begin{tikzpicture}
        \begin{axis}[
            globalstyle,
            legend pos=north east,
            xlabel={Number of Basis},
            xmin=4, xmax=14,
            ymin=0.0
        %xtick={3,...,14},
        ]

            \addplot+[mark=*, blue] coordinates {
                (4, 0.033566) (5, 0.012962) (6, 0.012962) (7, 0.006028) (8, 0.006028) (9, 0.000312) (10, 0.000312) (11, 0.000675) (12, 0.000083) (13, 0.000426) (14, 0.000171)
            };
            \addlegendentry{\PredErr}

            \addplot[very thick, dashed, red] coordinates {(4, 0.01785714) (14, 0.01785714)};
            \addlegendentry{Uni-Diff}
        \end{axis}
    \end{tikzpicture}
    \caption{
        Approximation quality as a function of the number of basis functions.
        Left: \Fdelta, the bound derived from Theorem~\ref{th:central-approx}.
        Right: \PredErr\, the absolute difference between the estimated predictive probability and the
        central predictive probability.
        The dotted line (Uni-Diff) indicates the corresponding error obtained using a uniform prior.
    }
    \label{fig:error_comparision}
\end{figure}


Next, we use the approximation to estimate the central predictive probability $C(X_7\mid S_6)$.
Here, \cDist denotes the marginal distribution induced by the central prior $\theta^{-\hf}(1 - \theta)^{-\hf}$.
We consider an observed sequence $S_6$ containing two successes and four failures.
Under \cDist, the predictive probability of success is $\frac{5}{14}$.
Figure~\ref{fig:error_comparision} reports the absolute difference between the estimated predictive probability and
the central predictive probability, denoted by \PredErr.
For comparison, the absolute difference of the predictive probability obtained from the
uniform prior is also shown as the dotted line (Uni-Diff).

The results show that as the number of basis functions increases, both \Fdelta\ and \PredErr\ decrease substantially.
This suggests that improvements in the approximation bound are accompanied by improvements in the resulting
predictive distribution.
Moreover, with a sufficient number of basis functions, the predictive estimate becomes considerably closer to the
central posterior than the estimate obtained from the uniform prior.

